\begin{lexlean}{Order}
\useglossary{lexlean.std.nat@1.0.0}
\useglossary{peano.arith@1.0.0}
\importmodule{Addition}
\title{Natural number order}

\begin{section}{positivity}
\heading{Positivity}
  \begin{predicatedefinition}{positive}{peano.arith::positive}
  \noaxioms
  For every natural number \(n\), \(n\) is positive holds exactly when \(0 < n\).
  \end{predicatedefinition}

  \begin{theorem}{successor-positive}
  \noaxioms
  For every natural number \(n\), the successor of \(n\) is positive.
  \begin{proof}
  Close the goal with \(zeroltsucc(n)\).
  \end{proof}
  \end{theorem}
\end{section}

\begin{section}{induction}
\heading{Order}
  \begin{theorem}{zero-le}
  \noaxioms
  For every natural number \(n\), \(0 ≤ n\).
  \begin{proof}
  \begin{induction}{n}
  \begin{case}{lexlean.std.nat::zero}
  \bind{}
  Close the goal with \(lerefl(0)\).
  \end{case}
  \begin{case}{lexlean.std.nat::succ}
  \bind{m;ih}
  Close the goal with \(letrans(ih, lesucc(m))\).
  \end{case}
  \end{induction}
  \end{proof}
  \end{theorem}

  \begin{theorem}{le-add-right}
  \noaxioms
  For every natural number \(n\), \(n ≤ n + 0\).
  \begin{proof}
  \begin{rewrite}{goal}
  \forward{\reference{Addition::add-zero}(n)}
  \end{rewrite}
  Close the goal with \(lerefl(n)\).
  \end{proof}
  \end{theorem}
\end{section}

\begin{section}{case-analysis}
\heading{Case analysis}
  \begin{theorem}{zero-or-successor}
  \noaxioms
  For every natural number \(n\), \(n = 0\) or there exists a natural number \(m\) such that \(n = succ(m)\).
  \begin{proof}
  \begin{cases}{n}
  \begin{case}{lexlean.std.nat::zero}
  \bind{}
  Select the left alternative.
  Close the goal by reflexivity.
  \end{case}
  \begin{case}{lexlean.std.nat::succ}
  \bind{k}
  Select the right alternative.
  Use \(k\) as the witness.
  Close the goal by reflexivity.
  \end{case}
  \end{cases}
  \end{proof}
  \end{theorem}

  \begin{theorem}{successor-not-zero}
  \allowaxioms{propext}
  For every natural number \(n\), not \(succ(n) = 0\).
  \begin{proof}
  Assume \(h\).
  Apply \(succnezero(n)\).
  Close the goal with \(h\).
  \end{proof}
  \end{theorem}
\end{section}
\end{lexlean}
